En este informe se pudo determinar diferencias entre la complejidad asintótica teórica y su comportamiento en la práctica, mediante la obtención de resultados replicables, y un análisis detallado.\\
Para los algoritmos de ordenamiento de arreglos, es posible confirmar la hipótesis de que el algoritmo más eficiente, es decir, el que se demora el menor tiempo para ejecutarse, es el algoritmo \texttt{Sort}, para todos los tipos de entrada. Y, para el algoritmo menos eficiente, se tiene el algoritmo \texttt{MergeSort}.\\
Para los algoritmos de multiplicación de matrices, el algoritmo \texttt{Naive} es consistentemente más eficiente para todos los tipos de matrices de entrada. Lo que significa que el algoritmo \texttt{Strassen} tiene el peor rendimiento de los dos.\\
En conclusión, con los resultados se evidencian diferencias significativas entre la complejidad asintótica teórica y su desempeño en la práctica. La sobrecarga asociada a la gestión de memoria y la recursividad puede contrarrestar las ventajas teóricas en escenarios reales, y solo un análisis experimental riguroso, con resultados replicables y un examen detallado, permite identificar la opción más eficiente en cada contexto.
